\documentclass[11pt]{article}

% PACKAGE IMPORTS
\usepackage{fancyhdr}
\usepackage[utf8]{inputenc}
\usepackage[T1]{fontenc}
\usepackage[english]{babel}
\usepackage{graphicx}
\usepackage{enumitem}
\usepackage{xcolor}
\usepackage{amsmath,amssymb}
\usepackage{booktabs}
\usepackage{tabularx}
\usepackage{geometry}
\usepackage{lastpage}
\usepackage[scaled]{helvet}
\usepackage{siunitx}
\sisetup{detect-all}
\PassOptionsToPackage{obeyspaces,spaces}{url}
\usepackage{hyperref}
\hypersetup{hidelinks}
\usepackage{url}
\usepackage{tikz} % Used for positioning text on title
\usepackage{eso-pic}
\usepackage[useregional]{datetime2}

\graphicspath{{pictures/template/}}

% FONT SETTINGS
\renewcommand\familydefault{\sfdefault}

% COLORS
\definecolor{ertred}{RGB}{226,12,24}

% MARGINS
\geometry{
a4paper,
inner=20mm,
outer=20mm,
top=26mm,
bottom=30mm,
footskip=10mm,
headsep=13mm
}

% LISTS AND PATHS
\setlist{itemsep=2pt,topsep=3pt}
\setlist[enumerate,1]{leftmargin=*}
\setlist[itemize]{leftmargin=*}
\urlstyle{tt}
\newcommand{\file}[1]{\path{#1}}

% FRONT PAGE BACKGROUND IMAGE
\newcommand\BackgroundPic{
\put(0,0){
\parbox[b][\paperheight]{\paperwidth}{
\vfill
\centering
\includegraphics[width=\paperwidth,height=\paperheight]{pictures/template/c_title-page.png} % Background title page image
\vfill
}
}
}

\begin{document}

\AddToShipoutPicture*{\BackgroundPic}
\ClearShipoutPicture

% HEADER AND FOOTER
\let\oldheadrule\headrule
\renewcommand{\headrule}{\color{ertred}\oldheadrule}
\renewcommand{\headrulewidth}{1.5pt}
\renewcommand{\footrule}{\hbox to\headwidth{\color{ertred}\leaders\hrule height \footrulewidth\hfill}}
\renewcommand{\footrulewidth}{1.5pt}
\setlength\headheight{48pt}

\pagestyle{fancy}
\fancyhead[C]{\rule[-2ex]{0pt}{2ex} \bf{\fontsize{12}{5}\selectfont WEISSHORN NOSE CONE -- FLUENT WORKFLOW}}
\lhead{\includegraphics[height=1.5cm]{ert_logo_small}\vspace{-4mm}}
\rhead{\includegraphics[height=1.5cm]{c_patch.png}\vspace{-4mm}}
\cfoot{}
\fancyfoot[R]{\bf{\thepage/\pageref{LastPage}}}
\fancyfoot[L]{\rule[-2ex]{0pt}{2ex}\bf{EPFL ROCKET TEAM}}

% TITLE PAGE
\begin{titlepage}
\vspace*{-2.40cm} % Push title to the top of the page
\noindent
\hspace*{-1.2cm}
\parbox{\linewidth}{%
\bfseries\fontsize{34}{38}\selectfont
FLUENT WORKFLOW -\\
NOSE CONE CFD\\
}
\newline
\hspace*{-1.2cm}
\parbox{\linewidth}{%
\bfseries\fontsize{14}{18}\selectfont
Fabio Feig - 315169 \\ % Names and scipers
}

% ADDITIONAL TITLE PAGE DETAILS
\begin{tikzpicture}[remember picture,overlay]
\node[anchor=south east,inner sep=39mm] at (current page.south east) {
\parbox[b][-1.6cm][t]{15cm}{
\raggedleft
\fontsize{12}{12} \selectfont \normalfont \textbf{Weisshorn -- Nose Cone CFD Workflow}
}
};
\end{tikzpicture}

\begin{tikzpicture}[remember picture,overlay]
\node[anchor=south east,inner sep=39mm] at (current page.south east) {
\parbox[b][-0.8cm][t]{15cm}{
\raggedleft
\fontsize{12}{12} \selectfont \normalfont Date: \today % Date
}
};
\end{tikzpicture}

\end{titlepage}

\noindent Weisshorn nose cone: Von K\'arm\'an, $L = \SI{810}{mm}$, $D = \SI{180}{mm}$, tip radius \SI{4.5}{mm}.
Verification runs of recommendation R6 of 2026\_C\_FD\_REPORT-1\_NOSE\_CONE.\par
\vspace{2mm}

% ---------------------------------------------------------------------------------
\section{Solver: Fluent}

Use \textbf{Fluent} (the ``Fluid Flow (Fluent)'' system in Workbench) with these settings: density-based,
implicit, steady, k-$\omega$ SST, energy on, air as an ideal gas with Sutherland viscosity. That is the
setup the nose cone report already prescribes (R6 and Section~8), and it is what the team runs on the
DLL VM (Fluent 2024 R2). Discovery is fine as the geometry editor inside Workbench, but do not use its
own fluid simulation for shocks.

\paragraph{Two models, not one.}
\begin{enumerate}
\item \textbf{2D axisymmetric} for the three zero-incidence runs ($M = 0.95$, 1.2, 1.47). It is cheap,
it resolves the boundary layer and the shocks properly, and it compares directly with the axisymmetric
Euler results.
\item \textbf{3D half-model} with a symmetry plane, used only for the \SI{5}{\degree} incidence run
($C_{N\alpha}$ and CP). The STEP file goes here.
\end{enumerate}

\paragraph{Input files.}
\begin{itemize}
\item Design solid (STEP): \file{nose_vonkarman_LD4.5_rn4.5mm_FINAL.step}, in \file{output/weisshorn/cad/}
of this repository. \file{output/} is not tracked by git; \texttt{python -m nosecone export configs/weisshorn.yaml}
writes it.\\
Use the STEP, not the STL: the STL opens as a faceted body (about 122\,600 triangles), on which
subtracting it from a flow enclosure, naming faces and meshing do not work cleanly. The STEP is a
true solid of revolution (nose, \SI{60}{mm} shoulder, base disc) in mm, with the $x$ axis as the
rocket axis and the tip at the origin. It is a filled solid with no wall thickness, which is what
external-flow CFD needs.
\item Profile points: \file{nose_profile_final.csv} (columns \texttt{x\_mm}, \texttt{r\_mm}), same folder.
\item Euler reference, in \file{data/}:
\file{euler_results.csv} and \file{cfd/cp/f45_vonkarman_f4.5_rn0.05_medium_M*.npz}
($M = 1.2$, 1.4, 1.6).
\end{itemize}

\paragraph{Free-stream conditions} from the nominal trajectory:

\begin{center}
\begin{tabular}{@{}cccccccc@{}}
\toprule
$M$ & altitude & $V$ & $p$ & $T$ & $\mathrm{Re}_L$ & first cell for $y^+ \approx 1$ \\
\midrule
0.95 & \SI{928}{m}  & \SI{321}{m/s} & \SI{90666}{Pa} & \SI{282.1}{K} & $1.7\cdot10^{7}$ & \SI{1.4}{\micro m} \\
1.20 & \SI{1478}{m} & \SI{402}{m/s} & \SI{84789}{Pa} & \SI{278.5}{K} & $2.0\cdot10^{7}$ & \SI{1.2}{\micro m} \\
1.47 & \SI{2189}{m} & \SI{488}{m/s} & \SI{77649}{Pa} & \SI{273.9}{K} & $2.3\cdot10^{7}$ & \SI{1.0}{\micro m} \\
\bottomrule
\end{tabular}
\end{center}

\noindent The first-cell heights are flat-plate estimates, so check $y^+$ after the first run. The
hot-day case reaches $M = 1.56$ (at \SI{2315}{m}) if you want one extra run.

% ---------------------------------------------------------------------------------
\clearpage
\section{2D workflow}

\begin{enumerate}
\item \textbf{Workbench:} add a ``Fluid Flow (Fluent)'' system. \textbf{Before you open the geometry},
set Geometry cell $\rightarrow$ Properties $\rightarrow$ Analysis Type = \textbf{2D}. If you skip this,
Fluent will not let you choose axisymmetric.

\item \textbf{Geometry:}
\begin{itemize}
\item Build one surface in the $XY$ plane with $y \geq 0$: a far-field rectangle minus the nose
half-section, with the $x$ axis as the rocket axis. The STEP already uses $x$ as the axis, and
\file{nose_profile_final.csv} gives the profile points.
\item Run the cylinder straight through to the outlet. With no base there is no wake, so you get
clean forebody numbers, the same as the Euler nose-plus-cylinder setup.
\item Domain size: about $20D$ upstream and radially (\SI{3.6}{m}) and about $10D$ behind the shoulder.
\item Split the wall at the shoulder and add named selections: \texttt{inlet}, \texttt{farfield},
\texttt{outlet}, \texttt{axis}, \texttt{wall\_nose}, \texttt{wall\_body}.
\end{itemize}

\item \textbf{Mesh} (Ansys Meshing):
\begin{itemize}
\item Inflation on both walls: \SI{1}{\micro m} first layer, growth 1.2, about 45 layers. That is
about \SI{18}{mm} in total, which covers the roughly \SI{10}{mm} boundary layer at the shoulder.
\item Refine the tip to about \SI{0.2}{mm} cells, since the tip radius is only \SI{4.5}{mm}.
\end{itemize}

\item \textbf{Fluent setup:}
\begin{itemize}
\item Launcher: 2D, double precision. General: Density-Based, Steady, Axisymmetric.
\item \textbf{Operating pressure = \SI{0}{Pa}}, so every pressure you enter is absolute.
\item Boundary conditions:
\begin{itemize}
\item \texttt{inlet} and \texttt{farfield}: pressure-far-field with $M$, $p$ and $T$ from the table,
flow direction $(1, 0)$.
\item \texttt{outlet}: pressure-outlet at $p_\infty$.
\item \texttt{axis}: axis.
\item Walls: no-slip, adiabatic.
\end{itemize}
\item Reference values: area $\pi \cdot 0.18^2/4 = \SI{0.02545}{m^2}$, length \SI{0.18}{m}. Fluent
integrates axisymmetric forces over the full \SI{360}{\degree}, so this gives the same $C_D$
definition as the Euler run.
\item Methods: Roe-FDS flux, least-squares gradients. Run first order for about 300--500 iterations,
then switch to second order. Start the Courant number at 1 and raise it to 5--10 once the run is stable.
\item Initialise with FMG or hybrid initialisation, starting from the far-field values.
\end{itemize}

\item \textbf{Convergence:} monitor the drag coefficient on \texttt{wall\_nose}. Fluent splits it into
pressure and viscous parts. The run is converged when $C_D$ changes by less than \SI{0.1}{\percent}
over about 500 iterations. Residuals alone do not tell you that.

\item \textbf{Mesh study:} at $M = 1.47$ only, run coarse, medium and fine meshes (about $\times1.5$
per direction each step), plus one or two pressure-gradient adaptions at the shock.

\item \textbf{Compare with the Euler results:}
\begin{itemize}
\item Fluent's \textbf{pressure} part should match $C_{D,\mathrm{w}}$. $M = 1.2$ is directly in
\file{euler_results.csv}; $M = 1.47$ falls between the 1.4 and 1.6 values.
\item Plot $C_p$ along the nose against \file{data/cfd/cp/f45_vonkarman_f4.5_rn0.05_medium_M*.npz}.
\item $M = 0.95$ replaces the cubic transonic model of the report.
\end{itemize}
\end{enumerate}

% ---------------------------------------------------------------------------------
\section{\texorpdfstring{3D run at \SI{5}{\degree} incidence}{3D run at 5 deg incidence}}

\begin{itemize}
\item Use a second Fluent system in 3D with a half domain cut on the $XY$ plane, and set that plane as a
symmetry boundary. Flow direction is $(\cos\SI{5}{\degree}, \sin\SI{5}{\degree}, 0)$.
\item For the mesh, Fluent Meshing (watertight workflow, poly-hexcore with prism layers) handles this
better than Ansys Meshing.
\item Normal force is dominated by pressure, so $y^+ \approx 30$--50 (first cell about
\SIrange{30}{50}{\micro m}) is acceptable here and keeps the cell count down.
\item Report the normal force and the moment about the tip.
\begin{itemize}
\item A half model gives half the force, so double it (or halve the reference area).
\item $C_{N\alpha} = C_N / 0.0873$; compare it with \SI{2}{\per\radian}.
\item $x_{\mathrm{CP}} = M_{\mathrm{tip}} / N$; compare it with \SI{405}{mm}.
\end{itemize}
\end{itemize}

% ---------------------------------------------------------------------------------
\section{Running Fluent from Claude Code on the DLL VM}

Claude Code cannot be installed onto the VM from the laptop, even while you are logged in: the Remote
Desktop window gives no shell on the VM. Run the install yourself in PowerShell on the VM. It takes about
10 minutes following 2026\_C\_FD\_REPORT-2\_MCP-FLUENT:
\begin{itemize}
\item \textbf{Section 3.4} to install and sign in (\texttt{claude}), make the \texttt{.claude} junction and
stage the copy on \texttt{Z:\textbackslash}.
\item \textbf{Sections 3.5--3.6} for the wheel cache and the repo clone. You need the repo URL from the
guide's author, because the guide only has \texttt{<project-git-url>}.
\item \textbf{Section 4.2} to install PyFluent and the MCP package.
\item \textbf{Section 6.2} to register the MCP server, then \textbf{Section 7.1} to connect to a running Fluent.
\end{itemize}

\paragraph{Things to know.}
\begin{itemize}
\item \textbf{\texttt{C:\textbackslash} may be wiped} between sessions, so redo Section 6.2 when that happens.
\item Write \texttt{.mcp.json} \textbf{without a BOM}.
\item Start the server with \texttt{server/start-server}, not the Web Server menu.
\item The copy on the VM is a separate installation. It starts without the context or the project notes
of the Claude Code sessions on the laptop.
\item The MCP server only drives the Fluent solver: reading a case, initialising, iterating, TUI and
Scheme commands, screenshots. Geometry and meshing stay manual. It is useful for the Mach sweep and the
mesh study, by changing the far-field values through TUI commands and re-running.
\end{itemize}

\end{document}
